\chapter{Network and Command-Line Client}\label{app:network}

\section{Ports}

\begin{table}[htbp]
\centering
\caption{Network ports of the board}\label{tab:ports}
\begin{tabularx}{\linewidth}{@{}P{0.16\linewidth}L@{}}
\toprule
\textbf{Port} & \textbf{Function} \\
\midrule
TCP 9100 & HP-GL, raw print port. One connection is one job. \\
TCP 9300 & Setup and status console. One client at a time. \\
UDP 9302 & Discovery. \\
\bottomrule
\end{tabularx}
\end{table}

These ports are different from the ports of ClearCNC (8888, 8889, 10040) and
ClearCoreROS (9200--9203).

\section{Discovery}

To find boards, a computer sends the text \code{CLEARHPGL_DISCOVER?} as a
UDP broadcast to port 9302. Each board replies with one line:

\begin{lstlisting}[language=console]
CLEARHPGL ClearHPGL_Firmware IP=10.168.41.5 HPGL=9100 CONSOLE=9300 FW=0.2.0 IDENT=7475A
\end{lstlisting}

The \menu{Discover} button of the desktop app and the \cmd{discover} command
of the client use this method. A broadcast does not go through a router, so
discovery finds only boards on the local subnet.

\section{The Command-Line Client}

\file{Firmware/host/clearhpgl_cli.py} is a client for the board. It needs
Python~3. For COM ports, it also needs the \code{pyserial} package.

\begin{table}[htbp]
\centering
\caption{Commands of the client}\label{tab:cli}
\begin{tabularx}{\linewidth}{@{}P{0.42\linewidth}L@{}}
\toprule
\textbf{Command} & \textbf{Function} \\
\midrule
\code{discover} & Find boards on the local network. \\
\code{--host} \textit{ip} \code{send} \textit{file} & Send an HP-GL file to TCP port 9100. \\
\code{--serial COM3 send} \textit{file} & Send an HP-GL file to a COM port with XON/XOFF, as a plotter driver does. \\
\code{--host} \textit{ip} \code{query "OI;OS;"} & Send output instructions and show the replies. \\
\code{--host} \textit{ip} \code{console} \textit{command} & Send one console command, for example \code{console status}. \\
\bottomrule
\end{tabularx}
\end{table}

Examples:

\begin{lstlisting}[language=console]
python Firmware/host/clearhpgl_cli.py discover
python Firmware/host/clearhpgl_cli.py --host 10.168.41.5 send drawing.plt
python Firmware/host/clearhpgl_cli.py --host 10.168.41.5 query "OI;OS;OA;"
python Firmware/host/clearhpgl_cli.py --host 10.168.41.5 console set draw_mm_s 80
\end{lstlisting}

\begin{note}
\notelabel The \code{console} command connects to port 9300. If the desktop
app is connected over Ethernet, it has the console, and the client cannot
connect. Close the app connection first.
\end{note}

\section{Test without a Board}

\file{Firmware/host/fake_plotter.py} operates as a board on the computer. Use
it to try the desktop app without hardware:

\begin{lstlisting}[language=console]
python Firmware/host/fake_plotter.py --bind 127.0.0.1
\end{lstlisting}

Then connect the app to 127.0.0.1.
